\section{引言}

\subsection{研究背景}

火星着陆是深空探测中约束最强的任务阶段之一。苏联 Mars 3 曾实现首次软着陆但很快失联；此后美国与中国完成了可持续运行的火星表面任务\cite{nasaNssdcMars3,govTianwen2021}。动力下降段的起始高度、速度和持续时间随任务设计而变，但共同目标是在有限推进剂、推力和状态约束下满足触地位置与速度要求\cite{braun2007mars}。在此阶段，减速与末端机动主要依赖反推发动机，因此制导算法的鲁棒性、燃料代价和计算时延都必须结合具体任务与硬件验证。

火星科学实验室（MSL）任务在末端下降采用天空起重机方案，以缆绳将"好奇号"火星车放至地面\cite{steltzner2013msl}。Mars 2020 的"毅力号"在 Jezero Crater 着陆；其 EDL 系统包含 Range Trigger、地形相对导航（TRN）和记录下降过程的相机系统\cite{chen2021mars2020}。MEDLI2 传感器套件为进入、下降和着陆环境测量提供了数据。中国于 2021 年 5 月完成天问一号火星着陆任务，祝融号部署在乌托邦平原\cite{govTianwen2021}。这些任务说明，导航、制导、执行器和环境模型必须作为完整 EDL 系统共同验证。

实际动力下降同时受重力、气动、执行器与导航误差影响。本项目采用三自由度平动点质量近似，状态为位置 $r$、速度 $v$ 与对数质量 $\ln m$，并施加终端位置和速度约束。手写基线将六台发动机的单台最小推力 $T_{\min}$ 和留出姿控余量后的单台上界 $T_2=0.8T_{\max}$ 汇总为有效推力包络；它不把发动机关机建成离散决策。该下界造成的非凸性是凸化文献所讨论的问题之一\cite{acikmese2007lossless}。实时重规划的周期、可行性和安全降级策略则取决于目标硬件与闭环场景，不能由本项目的主机基准预先假定。

\subsection{研究现状}

动力下降轨迹优化可采用间接法、直接法和凸优化方法。间接法借助庞特里亚金极大值原理将问题转化为两点边值问题；协态初值、开关结构和路径约束会使数值求解困难。直接法将连续问题转录为有限维非线性规划。伪谱 Legendre 离散可在较少节点上获得高精度近似\cite{el2009legendre}；直接配点则提供分段约束处理的通用途径\cite{benson2003direct}。其规模、收敛性和时延取决于转录方式、约束与求解器设置，是否适合嵌入式使用需要目标问题和硬件上的实测证据。

凸优化方法为部分受限动力下降问题提供可分析的近似或等价表述。Açıkmeşe 和 Ploen 研究了火星动力下降的凸规划表述\cite{acikmese2007lossless}；Blackmore 等人研究了特定非线性最优控制问题的控制约束无损凸化条件\cite{blackmore2012lossless}。Liu、Lu 和 Pan 的 2017 年论文综述了航空航天凸优化应用\cite{liu2017trajectory}。这些结果说明了 SOCP 与序列凸化的理论基础，但本项目的离散、线性化和物理可行性仍须由自身审计而非文献标题推断。

\subsection{现有工作的局限}

工程实现仍需分别处理建模、数值与证据问题。通用建模层便于原型验证，却可能引入建模、规范化和数据转换开销；直接输入求解器矩阵则需要可审查的装配与符号验证。不同求解器、模型版本和运行环境的数值输出只有连同物理残差才具有可比性。单精度与双精度的差异也必须在相同模型、缩放和容差下控制实验，不能由不同配置的单次结果推导硬件选择。

\subsection{本文贡献}

本文在凸优化制导的理论框架下，以火星动力下降轨迹优化为具体场景，给出以下可复核贡献。

（1）\textbf{原始手写稀疏矩阵资产。} \texttt{MarsLanding/MarsLanding.c} 是作者原始、独立、人工维护的 CRS 构造及 CRS--CCS 转换实现。等式约束矩阵含 733 个非零元，不等式约束矩阵含 403 个非零元；自动生成和 Python 路径只承担交叉验证职责，不能覆盖或替代该资产。

（2）\textbf{分组数值对比。} 三个 C ECOS 目标可直接构建运行；另保存三条 Python 对照记录，分别覆盖 ECOS、Clarabel 和 IPOPT。两组均在一位小数目标值上记录为 400.7 kg。该对比用于发现明显实现回归，不声称高精度逐位相等，也不据此证明物理模型正确或凸化无损。

（3）\textbf{物理修正的终端时间研究。} 在保留原始手写基线的同时，建立显式干重与区间内下滑锥约束的参数化 SOCP，并用冻结网格搜索、逐点审计和 ECOS/Clarabel 复算记录可行边界。该研究发现节点可行不等于区间可行，并报告网格敏感性而非宣称连续时间最优。

（4）\textbf{可追溯证据链。} 场景、逐样本记录、聚合结果、claim ledger 和论文图表在仓库中关联；历史 400.7 kg 路径仅作为数值 golden，物理审计、性能和精度结论均明确其测量范围。作为范围受限的主机记录，N150 上 setup 后的 1000 次求解 CPU 时间均值为 7.361 ms。

\subsection{论文组织}

第 2 章定义三自由度模型和物理约定；第 3 章说明凸化、离散化及其适用边界；第 4 章给出原始独立手写 CRS--CCS 资产和 ECOS 接口约定；第 5 章给出分组数值证据、物理修正自由时间研究和主机计时范围；第 6 章界定部署证据与待验证事项；第 7 章总结结论与下一步验证路线。
